%% =====================================================================
%%  Unit 5 — What's this?  (Les objets autour de moi)
%%  Pistes : 5.01–5.04 dialogues · 5.05–5.06 vocabulaire · 5.07–5.08 grammaire
%%           5.09 prononciation · 5.10–5.13 exercices (textes en annexe)
%% =====================================================================

\newcommand{\otIllusUnV}{\otIllusDuo{pens/104/86/15/0.10,keys/190/92/18/0.09,books/122/14/13/0.08,
  blue/150/8/11/0.08,red/196/22/12/0.07,boxes/100/50/12/0.07}{sam}{What are these?}{lea}{They're my books!}}

\unite{5}{What's this?}{Les objets autour de moi}{otIllusUnV}

\situation{Épisode 5}{Lundi, deuxième semaine. Léa découvre enfin son bureau au studio : Emma lui
fait le tour de ses nouvelles affaires… et de celles de Jack, qui traînent un peu partout. Le samedi
suivant, Sam vient l'aider à déballer les cartons arrivés de Lyon.}

\begin{objectifs}
\item nommer les objets du quotidien et dire de quelle couleur ils sont ;
\item demander ce qu'est un objet : \en{What's this?} \en{What are these?} ;
\item employer \en{a}, \en{an} et \en{the}, et mettre un nom au pluriel ;
\item montrer ce qui est proche ou loin avec \en{this}, \en{that}, \en{these}, \en{those} ;
\item prononcer le \en{-s} du pluriel : \phon{s}, \phon{z} ou \phon{ɪz}.
\end{objectifs}

\auprogramme
  {Nommer, montrer et décrire des objets}
  {\en{a}, \en{an}, \en{the} ; le pluriel\par\vspace{0.6mm}\en{this}, \en{that}, \en{these}, \en{those}}
  {Objets du quotidien, couleurs : 35 expressions}
  {Le pluriel : \phon{s}, \phon{z}, \phon{ɪz}}
  {Anglais américain et britannique}

\semaine{casque/otNavy/Dialogues/Écouter et lire les deux scènes,
  bulle/otGreenDark/Vocabulaire/Apprendre les 35 expressions,
  blocs/otViolet/Grammaire/Articles pluriel et this that,
  sons/otOrangeDark/Sons et culture/Le -s du pluriel et l'anglais américain,
  crayon/otNavy/Exercices/12 exercices corrigés,
  coche/otGreenDark/Bilan/L'essentiel et le test express}

\clearpage
%% ---------------------------------------------------------------------
\rubrique{ecouter}{Des objets et des cartons}

\begin{dialogue}{1}{Your new desk}{Lundi, 9 heures, au studio. Emma montre à Léa son nouveau bureau.}{5.01}{5.02}
\dit{emma}{\hl{This is} your desk, Léa.\note{1}}{Voici ton bureau, Léa.}
\dit{lea}{Great! \hl{What's this}?}{Génial ! C'est quoi, ça ?}
\dit{emma}{\hl{It's a} graphics tablet. It's new!}{C'est une tablette graphique. Elle est neuve !}
\dit{lea}{Cool! And \hl{what's that}?\note{2}}{Super ! Et ça, là-bas, c'est quoi ?}
\dit{emma}{\hl{That's the} coffee machine. It's very important!\note{3}}{C'est la machine à café. C'est très important !}
\dit{lea}{Ha ha! \hl{Are these} my \hl{pens}?\note{4}}{Ha ha ! Ce sont mes stylos ?}
\dit{emma}{Yes, \hl{they are}. And the \hl{blue notebook} is for you too.\note{5}}{Oui. Et le carnet bleu est pour toi aussi.}
\dit{lea}{Thank you! Is this your umbrella?}{Merci ! C'est ton parapluie ?}
\dit{emma}{No, it isn't. It's Jack's. And \hl{those} red keys are Jack's too!}{Non, c'est celui de Jack. Et ces clés rouges, là-bas, sont aussi à Jack !}
\dit{lea}{Jack and his things!}{Jack et ses affaires !}
\end{dialogue}

\begin{comprendre}
\pointnote{1} \en{\textbf{This is your desk.}} \en{this} montre ce qui est \textbf{près} de soi :
  « ceci », « ce… ci ». Avec \en{be}, \en{This is…} veut dire « Voici… ».
\pointnote{2} \en{\textbf{What's that?}} \en{that} montre ce qui est \textbf{plus loin} : « cela »,
  « ce… là ». \en{What's this?} = « Qu'est-ce que c'est ? » pour un objet proche.
\pointnote{3} \en{\textbf{the coffee machine}} \en{the} (« le, la, les ») désigne une chose précise,
  connue de tous : il n'y a qu'une machine à café au studio.
\pointnote{4} \en{\textbf{these pens}} \en{these} est le pluriel de \en{this}, et \en{those} celui
  de \en{that}. Le nom prend un \en{-s} au pluriel : \en{pen} → \en{pens}.
\pointnote{5} \en{\textbf{the blue notebook}} L'adjectif se place \textbf{avant} le nom et ne
  s'accorde jamais : \en{a blue pen}, \en{two blue pens}. \en{It's Jack's} = « c'est celui de Jack ».
\end{comprendre}

\begin{dialogue}{2}{What are these?}{Samedi, chez Léa, à Brooklyn. Sam l'aide à ouvrir les cartons
  arrivés de France.}{5.03}{5.04}
\dit{sam}{Wow, \hl{boxes}! What's in this box?}{Waouh, des cartons ! Qu'est-ce qu'il y a dans ce carton ?}
\dit{lea}{\hl{Books} and \hl{photos}.}{Des livres et des photos.}
\dit{sam}{And \hl{what are these}?}{Et ça, qu'est-ce que c'est ?}
\dit{lea}{\hl{They're} my \hl{glasses}\ldots{} and my old \hl{watches}!\note{6}}{Ce sont mes lunettes… et mes vieilles montres !}
\dit{sam}{Is that a lamp?}{C'est une lampe, ça ?}
\dit{lea}{Yes, it's \hl{an old lamp} from Lyon. It's orange!}{Oui, c'est une vieille lampe de Lyon. Elle est orange !}
\dit{sam}{Cool. And \hl{those plants}?}{Cool. Et ces plantes, là-bas ?}
\dit{lea}{They're from Mrs.\ Lopez. They're \hl{a present}!\note{7}}{Elles viennent de Mrs.\ Lopez. C'est un cadeau !}
\dit{sam}{She's so nice. And where's your TV?}{Elle est trop gentille. Et où est ta télé ?}
\dit{lea}{No TV for me! My laptop is my TV.}{Pas de télé pour moi ! Mon ordinateur portable me sert de télé.}
\end{dialogue}

\begin{comprendre}
\pointnote{6} \en{\textbf{glasses}, \textbf{watches}} Après \en{s}, \en{x}, \en{ch} ou \en{sh}, le
  pluriel s'écrit \en{-es} et se prononce \phon{ɪz} : \en{box} → \en{boxes}, \en{watch} →
  \en{watches}. \en{glasses} (« des lunettes ») est toujours au pluriel.
\pointnote{7} \en{\textbf{They're a present.}} « C'est un cadeau » (mot à mot : « elles sont un
  cadeau »). Quand on parle de plusieurs choses, on dit \en{They're…} et non \en{It's…}.
\end{comprendre}

%% ---------------------------------------------------------------------
\rubrique{mots}{Les mots de l'unité}

Voici les 35 expressions de l'unité. Les noms sont donnés avec \en{a} ou \en{an} : apprenez-les
toujours ainsi, l'article fait partie du mot.

\begin{vocabulaire}{Les objets du quotidien}{5.05}
\categorie{Au bureau}
\mot{a pen}{ə ˈpen}{un stylo}
\mot{a pencil}{ə ˈpensəl}{un crayon}
\mot{a notebook}{ə ˈnoʊtbʊk}{un carnet, un cahier}
\mot{a book}{ə ˈbʊk}{un livre}
\mot{a desk}{ə ˈdesk}{un bureau (le meuble)}
\mot{a chair}{ə ˈtʃer}{une chaise}
\mot{a laptop}{ə ˈlæptɑːp}{un ordinateur portable}
\mot{a phone}{ə ˈfoʊn}{un téléphone}
\categorie{Sur moi}
\mot{a bag}{ə ˈbæɡ}{un sac}
\mot{a key}{ə ˈkiː}{une clé}
\mot{a watch}{ə ˈwɑːtʃ}{une montre}
\mot{glasses}{ˈɡlæsɪz}{des lunettes}
\mot{an umbrella}{ən ʌmˈbrelə}{un parapluie}
\categorie{À la maison}
\mot{a box}{ə ˈbɑːks}{une boîte, un carton}
\mot{a lamp}{ə ˈlæmp}{une lampe}
\mot{a plant}{ə ˈplænt}{une plante}
\mot{a bottle}{ə ˈbɑːtəl}{une bouteille}
\mot{a cup}{ə ˈkʌp}{une tasse}
\mot{a TV}{ə ˌtiːˈviː}{une télé}
\mot{a present}{ə ˈprezənt}{un cadeau}
\end{vocabulaire}

\begin{vocabulaire}{Les couleurs et les descriptions}{5.06}
\categorie{Les couleurs}
\mot{red}{red}{rouge \otNuance{E0574C}}
\mot{blue}{bluː}{bleu \otNuance{2F6FD6}}
\mot{green}{ɡriːn}{vert \otNuance{3B9B4E}}
\mot{yellow}{ˈjeloʊ}{jaune \otNuance{F2C53D}}
\mot{orange}{ˈɔːrɪndʒ}{orange \otNuance{F0903B}}
\mot{pink}{pɪŋk}{rose \otNuance{F29BC2}}
\mot{purple}{ˈpɝːpəl}{violet \otNuance{7B4FB3}}
\mot{brown}{braʊn}{marron \otNuance{8A5A3B}}
\mot{black}{blæk}{noir \otNuance{1C1C1C}}
\mot{white}{waɪt}{blanc \otNuance{FFFFFF}}
\mot{gray}{ɡreɪ}{gris \otNuance{9AA0A6} (britannique : \en{grey})}
\categorie{Décrire et demander}
\mot{new}{nuː}{neuf, nouveau}
\mot{old}{oʊld}{vieux, ancien}
\mot{small}{smɑːl}{petit}
\mot{What color is it?}{ˌwʌt ˈkʌlɚ ɪz ɪt}{De quelle couleur est-il / elle ?}
\end{vocabulaire}

%% ---------------------------------------------------------------------
\rubrique{grammaire}{Les articles, le pluriel, this et that}

\otIntertitre[otViolet]{1 · \en{a}, \en{an} ou \en{the} ?}
\begin{paire}
\begin{regle}[\en{a} / \en{an} : un, une]
Un objet \textbf{parmi d'autres}, ou dont on parle pour la première fois :\par\vspace{1mm}
\exen{It's a pen.}{C'est un stylo.}
\exen{It's an umbrella.}{C'est un parapluie.}
\end{regle}
\begin{regle}[\en{the} : le, la, les]
Un objet \textbf{précis}, déjà connu ou unique ; \en{the} ne change ni au féminin, ni au pluriel :\par\vspace{1mm}
\exen{The pen is blue.}{Le stylo (dont on parle) est bleu.}
\exen{the coffee machine}{la machine à café (du studio)}
\end{regle}
\end{paire}

\begin{astuce}[La prononciation de \en{the}]
Devant un son consonne, \en{the} se dit \phon{ðə} : \en{the pen}, \en{the desk}. Devant un son
voyelle, on dit \phon{ði} : \en{the umbrella}, \en{the old lamp}. C'est la même règle que pour
\en{a} et \en{an}.
\end{astuce}

\otIntertitre[otViolet]{2 · Le pluriel des noms}
\begin{center}\otTableau
\begin{tabular}{@{}l@{\hspace{7mm}}>{\enOn}l<{\enOff}@{\hspace{7mm}}>{\enOn\bfseries}l<{\enOff}@{}}
\otTH{Règle} & \otTH{Singulier} & \otTH{Pluriel}\\\midrule
cas général : \en{+ s}                         & a pen, a key, a book        & pens, keys, books\\
après \en{s}, \en{x}, \en{ch}, \en{sh} : \en{+ es} & a box, a watch, a glass   & boxes, watches, glasses\\
consonne \en{+ y} → \en{ies}                    & a city, a baby              & cities, babies\\
\en{f} ou \en{fe} → \en{ves}                    & a wife, a knife             & wives, knives\\\midrule
pluriels irréguliers                             & a man, a woman, a child     & men, women, children\\
                                                 & a person                    & people\\\bottomrule
\end{tabular}\end{center}

\blocaudio{5.07}{Écoutez et répétez}{Six noms au singulier, puis au pluriel, dont les pluriels
irréguliers : \en{man}, \en{men} ; \en{child}, \en{children}…}

\otIntertitre[otViolet]{3 · \en{this}, \en{that}, \en{these}, \en{those}}
\begin{center}
\begin{tikzpicture}
  \node[fill=otVioletTint,rounded corners=1.6mm,minimum width=56mm,minimum height=8mm,font=\otSemi\fontsize{8}{9}\selectfont] at (35mm,24mm) {\otCaps[10]{près de moi}};
  \node[fill=otVioletTint,rounded corners=1.6mm,minimum width=56mm,minimum height=8mm,font=\otSemi\fontsize{8}{9}\selectfont] at (95mm,24mm) {\otCaps[10]{loin de moi}};
  \node[anchor=east,font=\otSemi\fontsize{8}{9}\selectfont,text=otViolet] at (5mm,12mm) {\otCaps[10]{singulier}};
  \node[anchor=east,font=\otSemi\fontsize{8}{9}\selectfont,text=otViolet] at (5mm,2mm) {\otCaps[10]{pluriel}};
  \node[font=\fontsize{11}{13}\selectfont] at (35mm,12mm) {\enb{\textbf{this} key} · cette clé-ci};
  \node[font=\fontsize{11}{13}\selectfont] at (95mm,12mm) {\enb{\textbf{that} key} · cette clé-là};
  \node[font=\fontsize{11}{13}\selectfont] at (35mm,2mm) {\enb{\textbf{these} keys} · ces clés-ci};
  \node[font=\fontsize{11}{13}\selectfont] at (95mm,2mm) {\enb{\textbf{those} keys} · ces clés-là};
\end{tikzpicture}
\end{center}

\blocaudio{5.08}{Écoutez et répétez}{Une question avec \en{this}, \en{that}, \en{these} ou \en{those},
puis sa réponse : \en{It's…} pour une chose, \en{They're…} pour plusieurs.}

\begin{attention}[Les pièges des francophones]
\fauxjuste{two blues pens}{two blue pens}
\fauxjuste{a umbrella}{an umbrella}
\fauxjuste{This keys are red.}{These keys are red.}
\fauxjuste{three childs}{three children}
\fauxjuste{What are this?}{What's this? / What are these?}
\end{attention}

%% ---------------------------------------------------------------------
\rubrique[Le -s du pluriel]{sons}{Le \en{-s} du pluriel : \phon{s}, \phon{z} ou \phon{ɪz}}

Le \en{-s} du pluriel se prononce toujours, mais de trois façons différentes, selon le son qui le
précède. Inutile d'apprendre une liste : votre bouche choisit d'elle-même le son le plus facile.

\begin{tcbraster}[raster columns=3,raster equal height=rows,raster column skip=3mm,raster before skip=2mm,raster after skip=3mm]
\otCartePrg{sons}{otOrangeDark}{\phon{s}}{après \phon{p}, \phon{t}, \phon{k}, \phon{f} (sons « sans voix ») : \en{books}, \en{cups}, \en{desks}}
\otCartePrg{sons}{otOrangeDark}{\phon{z}}{après une voyelle ou un son « avec voix » : \en{pens}, \en{bags}, \en{keys}, \en{phones}}
\otCartePrg{sons}{otOrangeDark}{\phon{ɪz}}{après \phon{s}, \phon{z}, \phon{ʃ}, \phon{tʃ}, \phon{dʒ} : une syllabe en plus ! \en{boxes}, \en{watches}}
\end{tcbraster}

\begin{paires}{5.09}{Écoutez le singulier, puis le pluriel, et répétez. Attention : \en{boxes} et
\en{watches} ont une syllabe de plus que \en{box} et \en{watch}.}
\sonpaire{book}{bʊk}{\phon{s}}{book\otH{s}}{bʊks}{des livres}
\sonpaire{map}{mæp}{\phon{s}}{map\otH{s}}{mæps}{des cartes, des plans}
\sonpaire{desk}{desk}{\phon{s}}{desk\otH{s}}{desks}{des bureaux}
\sonpaire{pen}{pen}{\phon{z}}{pen\otH{s}}{penz}{des stylos}
\sonpaire{bag}{bæɡ}{\phon{z}}{bag\otH{s}}{bæɡz}{des sacs}
\sonpaire{key}{kiː}{\phon{z}}{key\otH{s}}{kiːz}{des clés}
\sonpaire{box}{bɑːks}{\phon{ɪz}}{box\otH{es}}{ˈbɑːksɪz}{des boîtes}
\sonpaire{watch}{wɑːtʃ}{\phon{ɪz}}{watch\otH{es}}{ˈwɑːtʃɪz}{des montres}
\end{paires}

%% ---------------------------------------------------------------------
\rubrique{culture}{Anglais américain et anglais britannique}

\begin{culture}[Deux pays, deux vocabulaires]
Américains et Britanniques se comprennent sans peine, mais certains objets du quotidien n'ont pas
le même nom. Ce livre emploie les mots américains ; voici les principaux équivalents britanniques.\par\vspace{2mm}
\begin{tabularx}{\linewidth}{@{}>{\enOn\bfseries}l<{\enOff}@{\hspace{3mm}}>{\enOn}l<{\enOff}@{\hspace{3mm}}X@{\hspace{5mm}}>{\enOn\bfseries}l<{\enOff}@{\hspace{3mm}}>{\enOn}l<{\enOff}@{\hspace{3mm}}X@{}}
\otTH{US} & \otTH{UK} & \otTH{Français} & \otTH{US} & \otTH{UK} & \otTH{Français}\\
cell phone & mobile      & portable        & cookie    & biscuit   & biscuit\\
apartment  & flat        & appartement     & pants     & trousers  & pantalon\\
elevator   & lift        & ascenseur       & vacation  & holiday   & vacances\\
subway     & underground & métro           & trash can & bin       & poubelle\\
fall       & autumn      & automne         & check     & bill      & addition\\
\end{tabularx}
\end{culture}

\begin{culture}[Et l'orthographe ?]
Quelques mots s'écrivent différemment : \en{color} / \en{colour}, \en{gray} / \en{grey},
\en{center} / \en{centre}, \en{neighbor} / \en{neighbour}. Les deux orthographes sont justes : l'important
est de rester cohérent dans un même texte.
\end{culture}

%% ---------------------------------------------------------------------
\rubrique{pratique}{Exercices}

\exercice[ecrire,niveau=1]{Au pluriel}{Écrivez le pluriel de chaque nom.}
\begin{questions}[2]
\item a pen → two \trou[26mm]{pens}
\item a box → two \trou[26mm]{boxes}
\item a watch → two \trou[26mm]{watches}
\item a city → two \trou[26mm]{cities}
\item a key → two \trou[26mm]{keys}
\item a child → two \trou[26mm]{children}
\item a man → two \trou[26mm]{men}
\item a woman → two \trou[26mm]{women}
\end{questions}

\exercice[ecrire,niveau=2]{\en{a}, \en{an} ou \en{the} ?}{Complétez.}
\begin{questions}
\item Look! It's \trou[12mm]{a} cat. \trou[12mm]{The} cat is black.
\item Is it \trou[12mm]{an} umbrella? Yes, and \trou[12mm]{the} umbrella is Jack's.
\item This is \trou[12mm]{an} old photo. \trou[12mm]{The} photo is from Lyon.
\item Where's \trou[12mm]{the} coffee machine, please?
\item London is \trou[12mm]{a} big city.
\end{questions}

\exercice[ecrire,niveau=1]{\en{this}, \en{that}, \en{these} ou \en{those} ?}{Complétez. L'indication
entre parenthèses dit si l'objet est près (ici) ou loin (là-bas).}
\begin{questions}[2]
\item \trou[16mm]{This} pen is blue. \textit{(ici)}
\item \trou[16mm]{Those} keys are Jack's. \textit{(là-bas)}
\item \trou[16mm]{These} books are new. \textit{(ici)}
\item \trou[16mm]{That} lamp is old. \textit{(là-bas)}
\item Is \trou[16mm]{this} your bag? \textit{(ici)}
\item Are \trou[16mm]{those} your plants? \textit{(là-bas)}
\end{questions}

\exercice[lire,ecrire,niveau=1]{Qu'est-ce que c'est ?}{Répondez par \en{It's a…}, \en{It's an…}
ou \en{They're…}.}
\begin{tcbraster}[raster columns=3,raster equal height=rows,raster column skip=3mm,raster row skip=3mm,raster before skip=1mm]
\otCarteObjet{objcle}{otNavy}{3}{They're keys.}
\otCarteObjet{objparapluie}{otRed}{1}{It's an umbrella.}
\otCarteObjet{objlivre}{otGreenDark}{2}{They're books.}
\otCarteObjet{objlampe}{otOrange}{1}{It's a lamp.}
\otCarteObjet{objmontre}{otViolet}{2}{They're watches.}
\otCarteObjet{objtasse}{otTeal}{1}{It's a cup.}
\end{tcbraster}

\exercice[ecrire,niveau=1]{Les couleurs}{Complétez avec une couleur.}
\begin{anglais}
\begin{questions}[2]
\item The sky is \trou[22mm]{blue}.
\item Snow is \trou[22mm]{white}.
\item Bananas are \trou[22mm]{yellow}.
\item Tomatoes are \trou[22mm]{red}.
\item Chocolate is \trou[22mm]{brown}.
\item Grass is \trou[22mm]{green}.
\end{questions}
\par\vspace{1mm}{\fontsize{8.6}{10}\selectfont\color{otMuted}\foreignlanguage{french}{\textit{the sky} : le ciel · \textit{snow} : la neige · \textit{grass} : l'herbe}}
\end{anglais}

\exercice[ecrire,niveau=2]{Corrigez l'erreur}{Chaque phrase contient une erreur. Récrivez-la correctement.}
\begin{questions}
\item I have two blues pens. \hfill \trou[70mm]{I have two blue pens.}
\item It's a umbrella. \hfill \trou[70mm]{It's an umbrella.}
\item These book are new. \hfill \trou[70mm]{These books are new.}
\item This keys are red. \hfill \trou[70mm]{These keys are red.}
\item Those is my bag. \hfill \trou[70mm]{That's my bag.}
\item The three childs are in London. \hfill \trou[70mm]{The three children are in London.}
\end{questions}

\exercice[ecouter,niveau=1,piste=5.10]{\phon{s}, \phon{z} ou \phon{ɪz} ?}{Écoutez chaque pluriel, dit deux
fois. Cochez le son de la terminaison.}
\begin{questions}[3]
\item cups \hfill \choix{\phon{s}} \choix{\phon{z}} \choix{\phon{ɪz}}
\item bags \hfill \choix{\phon{s}} \choix{\phon{z}} \choix{\phon{ɪz}}
\item boxes \hfill \choix{\phon{s}} \choix{\phon{z}} \choix{\phon{ɪz}}
\item desks \hfill \choix{\phon{s}} \choix{\phon{z}} \choix{\phon{ɪz}}
\item phones \hfill \choix{\phon{s}} \choix{\phon{z}} \choix{\phon{ɪz}}
\item watches \hfill \choix{\phon{s}} \choix{\phon{z}} \choix{\phon{ɪz}}
\item books \hfill \choix{\phon{s}} \choix{\phon{z}} \choix{\phon{ɪz}}
\item keys \hfill \choix{\phon{s}} \choix{\phon{z}} \choix{\phon{ɪz}}
\item glasses \hfill \choix{\phon{s}} \choix{\phon{z}} \choix{\phon{ɪz}}
\end{questions}

\exercice[ecouter,ecrire,niveau=2,piste=5.11]{C'est à qui ?}{Quatre petits dialogues. Écoutez, puis
complétez le tableau : à qui est l'objet, et de quelle couleur est-il ?}
\begin{center}\otTableau
\begin{tabular}{@{}l@{\hspace{12mm}}l@{\hspace{12mm}}l@{}}
\otTH{Objet} & \otTH{À qui ?} & \otTH{Couleur}\\\midrule
{\otSemi 1 · \enb{the phone}}    & \trou[34mm]{Emma}        & \enb{\trou[22mm]{pink}}\\[1mm]
{\otSemi 2 · \enb{the keys}}     & \trou[34mm]{Léa}         & \enb{\trou[22mm]{green}}\\[1mm]
{\otSemi 3 · \enb{the umbrella}} & \trou[34mm]{Mr. Carter}  & \enb{\trou[22mm]{black}}\\[1mm]
{\otSemi 4 · \enb{the bag}}      & \trou[34mm]{Maya}        & \enb{\trou[22mm]{purple}}\\\bottomrule
\end{tabular}\end{center}

\exercice[ecouter,ecrire,niveau=2,piste=5.12]{Dictée à trous}{Sam et Léa ouvrent les derniers cartons.
Écoutez et complétez.}
\begin{anglais}
\noindent\begin{tabular}{@{}l@{\hspace{6mm}}l@{}}
\otNomDit{sam} & What's \trou{this}?\\[1.8mm]
\otNomDit{lea} & It's \trou{a} lamp from Lyon.\\[1.8mm]
\otNomDit{sam} & And what are \trou{these}?\\[1.8mm]
\otNomDit{lea} & They're my \trou{books}.\\[1.8mm]
\otNomDit{sam} & Are \trou{those} your plants?\\[1.8mm]
\otNomDit{lea} & Yes, \trou{they} are. They're from Mrs.\ Lopez.\\[1.8mm]
\otNomDit{sam} & \trou{The} lamp is orange!\\[1.8mm]
\otNomDit{lea} & Yes, and all my \trou{boxes} are brown!\\
\end{tabular}
\end{anglais}

\exercice[ecrire,niveau=2]{Traduisez}{Traduisez en anglais.}
\begin{questions}[1][french]
\item Qu'est-ce que c'est ? — C'est un parapluie.\lignes{1}
\item Ce sont mes clés.\lignes{1}
\item Ce livre-ci est bleu.\lignes{1}
\item Ces stylos-là sont à Jack.\lignes{1}
\item Les enfants sont à Londres.\lignes{1}
\item C'est un vieux sac noir.\lignes{1}
\item De quelle couleur est ta lampe ? — Elle est orange.\lignes{1}
\item Ces boîtes sont petites.\lignes{1}
\end{questions}

\exercice[ecrire,niveau=3]{Sur votre bureau}{Décrivez cinq objets de votre bureau ou de votre sac,
avec leur couleur. Utilisez \en{This is…}, \en{These are…}, \en{It's…}, \en{They're…}.}
\lignes{5}

\exercice[ecouter,parler,niveau=3,piste=5.13]{À vous de parler !}{Emma vous aide à ranger le studio.
Répondez-lui \textbf{à voix haute} pendant chaque pause.}
\begin{center}\otTableau
\begin{tabular}{@{}>{\otSemi\fontsize{8.6}{10}\selectfont\color{otNavy}}l@{\hspace{3mm}}>{\enOn}p{70mm}<{\enOff}@{\hspace{6mm}}p{78mm}@{}}
 & \otTH{Emma dit\ldots} & \otTH{Vous\ldots}\\\midrule
1. & Is this your umbrella? & répondez non, et dites que c'est celui de Jack.\\
2. & What color is your bag? & donnez la couleur de votre sac.\\
3. & What are these? & dites que ce sont vos clés.\\
4. & Is that your laptop? & répondez oui, par une réponse courte.\\
5. & Oh, it's new! & répondez que non, il est vieux.\\\bottomrule
\end{tabular}\end{center}

%% ---------------------------------------------------------------------
\rubrique{bilan}{Faites le point}

\begin{essentiel}
\begin{minipage}[t]{0.52\linewidth}
\otEssTitre{Les phrases clés}
\renewcommand{\arraystretch}{1.32}
\begin{tabular}{@{}>{\enOn\otSemi\fontsize{9.4}{11}\selectfont}l<{\enOff}@{\hspace{3mm}}>{\fontsize{9.2}{11}\selectfont\color{white!78!otNavy}}l@{}}
What's this? — It's a pen.   & C'est un stylo.\\
What are these?              & Qu'est-ce que c'est ? (plur.)\\
They're my keys.             & Ce sont mes clés.\\
What color is it?            & De quelle couleur est-il ?\\
It's an old blue bag.        & C'est un vieux sac bleu.\\
It's Jack's.                 & C'est celui de Jack.\\
\end{tabular}
\end{minipage}\hfill
\begin{minipage}[t]{0.42\linewidth}
\otEssTitre{Articles}
{\fontsize{9.2}{11.6}\selectfont \en{a pen} (un) · \en{the pen} (le stylo précis)\par}
\otEssTitre{Pluriel}
{\fontsize{9.2}{11.6}\selectfont \en{+s} · \en{+es} · \en{-ies} · \en{men}, \en{women}, \en{children}, \en{people}\par}
\otEssTitre{Montrer}
{\enb{\fontsize{9.2}{11.6}\selectfont this / these (ici) · that / those (là-bas)}\par}
\otEssTitre{Prononciation}
{\fontsize{9.2}{11.6}\selectfont \en{books} \phon{s} · \en{pens} \phon{z} · \en{boxes} \phon{ɪz}\par}
\end{minipage}
\end{essentiel}

\begin{jesais}
\item nommer une vingtaine d'objets du quotidien et onze couleurs ;
\item demander ce qu'est un objet et répondre : \en{It's a…} / \en{They're…} ;
\item choisir entre \en{a}, \en{an} et \en{the} ;
\item mettre un nom au pluriel, y compris \en{men}, \en{women}, \en{children}, \en{people} ;
\item utiliser \en{this}, \en{that}, \en{these}, \en{those} ;
\item prononcer le pluriel \phon{s}, \phon{z} ou \phon{ɪz}.
\end{jesais}

\clearpage
\begin{testexpress}
\qcm[a]{« une clé » se dit :}{\en{a key}}{\en{an key}}{\en{the key}}
\qcm[b]{Le pluriel de \en{box} est :}{\en{boxs}}{\en{boxes}}{\en{boxies}}
\qcm[c]{Le pluriel de \en{child} est :}{\en{childs}}{\en{childrens}}{\en{children}}
\qcm[b]{« ces livres(-ci) » se dit :}{\en{this books}}{\en{these books}}{\en{those book}}
\qcm[a]{« cette lampe(-là) » se dit :}{\en{that lamp}}{\en{those lamp}}{\en{this lamp}}
\qcm[c]{« deux stylos bleus » se dit :}{\en{two blues pens}}{\en{two pens blue}}{\en{two blue pens}}
\qcm[b]{\en{What are these?} — }{\en{It's a pen.}}{\en{They're pens.}}{\en{They're a pen.}}
\qcm[c]{Dans \en{watches}, la terminaison se prononce :}{\phon{s}}{\phon{z}}{\phon{ɪz}}
\qcm[a]{« un téléphone portable » se dit, aux États-Unis :}{\en{a cell phone}}{\en{a mobile}}{\en{a handy}}
\qcm[a]{\en{the} se prononce \phon{ði} devant :}{un son voyelle}{un son consonne}{un nom au pluriel}
\end{testexpress}

\noindent{\fontsize{9.6}{12}\selectfont\color{otMuted}Niveau visé (CECRL, niveau A1) : « peut
nommer et décrire simplement des objets familiers ». Les corrigés de l'unité sont
page~\pageref{cor-u5-1} ; les textes des pistes d'exercices, page~\pageref{trans-5.10}.\par}
